% =====================================================================
%  مقدمة الناظم (٣): الأبيات ٤٤–٦٦ — الرموز والاصطلاحات
%  (أهم جزء في المقدمة: عليه يُبنى فهم بقية القصيدة)
% =====================================================================
\chapter{مقدمة الناظم (٣): الرموز والاصطلاحات}

\begin{tanbih}[لماذا هذا القسم مهم؟]
الشاطبي لم يذكر أسماء القراء صراحة في أكثر القصيدة، بل وضع لكل قارئ \textbf{حرفًا} يرمز إليه، ولبعض الجماعات \textbf{حرفًا أو كلمة}، ووضع قواعد يُفهم منها قراءة من لم يُذكر. فمن لم يفهم هذه الأبيات لم يفهم القصيدة.
\end{tanbih}

\section{الرموز الحرفية لكل قارئ وراوٍ}

\bayt{44}{وَهَا أَنَا ذَا أَسْعَى لَعَلَّ حُرُوفَهُمْ}{يَطُوعُ بِهَا نَظْمُ الْقَوَافِي مُسَهَّلَا}
\begin{mufradat}
\kalima{حروفهم}{قراءاتهم (الحرف: القراءة).}
\kalima{يطوع}{ينقاد.}
\end{mufradat}
\maana أسعى راجيًا أن تنقاد قراءاتهم لنظم القصيدة بسهولة.

\bayt{45}{جَعَلْتُ أَبَا جَادٍ عَلَى كُلِّ قَارِئٍ}{دَلِيلًا عَلَى الْمَنْظُومِ أَوَّلَ أَوَّلَا}
\begin{mufradat}
\kalima{أبا جاد}{حروف «أبجد هوز حطي كلمن سعفص قرشت».}
\kalima{أوّلَ أوّلا}{بالترتيب: الأول للأول، والذي بعده لمن بعده.}
\end{mufradat}
\maana جعلت حروف «أبجد» رموزًا للقراء والرواة بالترتيب الذي ذكرتهم به: الحرف الأول للقارئ الأول، والثاني لراويه الأول، والثالث لراويه الثاني، وهكذا.

\begin{qaida}[جدول الرموز الحرفية]
\centering
\renewcommand{\arraystretch}{1.5}
\begin{tabular}{c c c}
\toprule
\textbf{الإمام} & \textbf{الراوي الأول} & \textbf{الراوي الثاني}\\
\midrule
\rzl{Nafi}{أ} \Nafi & \rzl{Nafi}{ب} \Qalun & \rzl{Nafi}{ج} \Warsh\\
\rzl{IbnKathir}{د} \IbnKathir & \rzl{IbnKathir}{هـ} \Bazzi & \rzl{IbnKathir}{ز} \Qunbul\\
\rzl{AbuAmr}{ح} \AbuAmr & \rzl{AbuAmr}{ط} \qname{AbuAmr}{الدوري} & \rzl{AbuAmr}{ي} \Susi\\
\rzl{IbnAmir}{ك} \IbnAmir & \rzl{IbnAmir}{ل} \Hisham & \rzl{IbnAmir}{م} \IbnDhakwan\\
\rzl{Asim}{ن} \Asim & \rzl{Asim}{ص} \Shuba & \rzl{Asim}{ع} \Hafs\\
\rzl{Hamza}{ف} \Hamza & \rzl{Hamza}{ض} \Khalaf & \rzl{Hamza}{ق} \Khallad\\
\rzl{Kisai}{ر} \Kisai & \rzl{Kisai}{س} \AbulHarith & \rzl{Kisai}{ت} \qname{Kisai}{الدوري}\\
\bottomrule
\end{tabular}

\smallskip
\raggedleft
\textbf{تذكّر}: «أبج — دهز — حطي — كلم — نصع — فضق — رست».
\end{qaida}

\bayt{46}{وَمِنْ بَعْدِ ذِكْرِي الْحَرْفَ أُسْمِي رِجَالَهُ}{مَتَى تَنْقَضِي آتِيكَ بِالْوَاوِ فَيْصَلَا}
\begin{mufradat}
\kalima{الحرف}{الكلمة القرآنية المختلف فيها.}
\kalima{أُسْمي رجاله}{أذكر رموز قرائه.}
\kalima{فيصلًا}{فاصلًا.}
\end{mufradat}
\maana عندي قاعدتان: (١) أذكر الكلمة القرآنية \textbf{أولًا} ثم أذكر رمز من قرأها \textbf{بعدها}. (٢) إذا انتهت المسألة أتيت بـ\textbf{واو} تفصلها عن المسألة التالية.

\begin{amthila}[مثال تطبيقي من البيت ١٠٨]
\begin{center}
\qw{وَمَالِكِ} يَوْمِ الدِّينِ \rz{Kisai}{رَاوِيهِ} \rz{Asim}{نَاصِرٌ} \quad \textbf{وَ}عِنْدَ سِرَاطٍ ...
\end{center}
\begin{itemize}
\item الكلمة القرآنية: «مالك» \ila{} ثم الرمز بعدها: \rz{Kisai}{ر} (\Kisai) و\rz{Asim}{ن} (\Asim).
\item ثم جاءت \textbf{الواو} في «وعند» فعرفنا أن المسألة انتهت وبدأت مسألة جديدة.
\end{itemize}
\end{amthila}

\bayt{47}{سِوَى أَحْرُفٍ لَا رِيبَةٌ فِي اتِّصَالِهَا}{وَبِاللَّفْظِ أَسْتَغْنِي عَنِ الْقَيْدِ إِنْ جَلَا}
\begin{mufradat}
\kalima{لا ريبة}{لا شك.}
\kalima{القيد}{الوصف المبيِّن للقراءة (كقوله: بالمد، بالفتح...).}
\kalima{جلا}{ظهر واتضح.}
\end{mufradat}
\maana
\textbf{الأول}: أنه قد يترك واو الفصل في مواضع قليلة لا يشك أحد أن الكلام فيها متصل بما قبله أو منفصل عنه.
\textbf{الثاني}: إذا كان نطقه بالكلمة يوضح القراءة فإنه لا يقيّدها بوصف.

\begin{amthila}
\begin{itemize}
\item \textbf{الاستغناء باللفظ}: قال «ومالك» بالألف، فعرفنا أن \Kisai{} و\Asim{} يقرآن بالألف دون أن يقول «بالمد».
\item \textbf{حرف لا ريبة فيه}: قال في البيت ١١٥: «... بالضم شمللا \textbf{كَمَا} بهم الأسباب»، فالكاف في «كما» ليست رمزًا لابن عامر، لأنها واضحة أنها كاف التشبيه قبل مثال.
\end{itemize}
\end{amthila}

\bayt{48}{وَرُبَّ مَكَانٍ كَرَّرَ الْحَرْفَ قَبْلَهَا}{لِمَا عَارِضٍ وَالْأَمْرُ لَيْسَ مُهَوِّلَا}
\begin{mufradat}
\kalima{لما عارضٍ}{لسبب عارض، كضرورة الوزن.}
\kalima{ليس مهوِّلًا}{ليس مخيفًا ولا صعبًا.}
\end{mufradat}
\maana قد يكرر الرمز (يأتي بالرمز مرتين للقارئ نفسه) قبل واو الفصل لضرورة الوزن أو لغيرها، وهذا لا يضر ولا يُشكل؛ لأن المعنى واحد.

\section{رموز الجماعة الحرفية: ث خ ذ ظ غ ش}

\begin{tanbih}
بعد «أبجد ... قرشت» بقيت ستة أحرف: «ثَخِذْ ظَغِشْ». جعل كل حرف منها رمزًا لجماعة من القراء.
\end{tanbih}

\bayt{49}{وَمِنْهُنَّ لِلْكُوفِي \rzg{ثَاءٌ} مُثَلَّثٌ}{وَسِتَّتُهُمْ \rzg{بِالْخَاءِ} لَيْسَ بِأَغْفَلَا}
\begin{mufradat}
\kalima{ثاء مثلث}{الثاء ذات النقط الثلاث.}
\kalima{ليس بأغفلا}{ليس خاليًا من النقط (الأغفل: الحرف غير المنقوط)، أي الخاء المعجمة لا الحاء.}
\end{mufradat}

\bayt{50}{عَنَيْتُ الْأُولَى أَثْبَتُّهُمْ بَعْدَ نَافِعٍ}{وَكُوفٍ وَشَامٍ \rzg{ذَالُهُمْ} لَيْسَ مُغْفَلَا}
\begin{mufradat}
\kalima{عنيت}{قصدت.}
\kalima{الأُولى أثبتّهم بعد نافع}{أي القراء الستة الذين ذكرتهم بعد نافع.}
\end{mufradat}

\bayt{51}{وَكُوفٍ مَعَ الْمَكِّيِّ \rzg{بِالظَّاءِ} مُعْجَمًا}{وَكُوفٍ وَبَصْرٍ \rzg{غَيْنُهُمْ} لَيْسَ مُهْمَلَا}
\begin{mufradat}
\kalima{معجمًا / ليس مهملًا}{منقوطًا. (المعجم: المنقوط، والمهمل: غير المنقوط.)}
\end{mufradat}

\bayt{52}{وَذُو النَّقْطِ \rzg{شِينٌ} لِلْكِسَائِي وَحَمْزَةٍ}{وَقُلْ فِيهِمَا مَعْ شُعْبَةٍ \rzg{صُحْبَةٌ} تَلَا}

\maana[(الأبيات ٤٩–٥٢)]
\begin{itemize}
\item \rzlg{ث} = الكوفيون الثلاثة: \Asim{} و\Hamza{} و\Kisai.
\item \rzlg{خ} = القراء الستة غير \Nafi: \IbnKathir، \AbuAmr، \IbnAmir، \Asim، \Hamza، \Kisai.
\item \rzlg{ذ} = الكوفيون و\IbnAmir{} الشامي.
\item \rzlg{ظ} = الكوفيون و\IbnKathir{} المكي.
\item \rzlg{غ} = الكوفيون و\AbuAmr{} البصري.
\item \rzlg{ش} = \Hamza{} و\Kisai.
\end{itemize}
ثم بدأ بالرموز الكلمية: فقال إن كلمة \rzg{صُحْبة} رمز لـ\Hamza{} و\Kisai{} ومعهما \Shuba.
وأكّد في الأبيات أن هذه الحروف \textbf{منقوطة} (الخاء لا الحاء، والذال لا الدال، والظاء لا الطاء، والغين لا العين، والشين لا السين) حتى لا تلتبس بحروف «أبجد».

\section{رموز الجماعة الكلمية}

\bayt{53}{\rzg{صِحَابٌ} هُمَا مَعْ حَفْصِهِمْ \rzg{عَمَّ} نَافِعٌ}{وَشَامٍ \rzg{سَمَا} فِي نَافِعٍ وَفَتَى الْعَلَا}
\bayt{54}{وَمَكٍّ وَ\rzg{حَقٌّ} فِيهِ وَابْنِ الْعَلَاءِ قُلْ}{وَقُلْ فِيهِمَا وَالْيَحْصَبِي \rzg{نَفَرٌ} حَلَا}
\bayt{55}{\rzg{وَحِرْمِيٌّ} الْمَكِّيُّ فِيهِ وَنَافِعٌ}{\rzg{وَحِصْنٌ} عَنِ الْكُوفِي وَنَافِعِهِمْ عَلَا}
\begin{mufradat}
\kalima{فتى العلا / ابن العلاء}{\AbuAmr.}
\kalima{اليحصبي}{\IbnAmir.}
\kalima{حرمي}{نسبة إلى الحَرَمين: مكة والمدينة.}
\end{mufradat}

\maana[(الأبيات ٥٢–٥٥)] هذه ثماني كلمات، كل كلمة رمز لجماعة:

\begin{qaida}[الرموز الكلمية الثمانية]
\centering
\renewcommand{\arraystretch}{1.5}
\begin{tabular}{c l l}
\toprule
\textbf{الكلمة} & \textbf{القراء} & \textbf{طريقة الحفظ}\\
\midrule
\rzlg{صُحْبَة} & \Hamza، \Kisai، \Shuba & «صحبة» فيها «شعبة»\\
\rzlg{صِحَاب} & \Hamza، \Kisai، \Hafs & «صحاب» مع «حفص»\\
\rzlg{عَمّ} & \Nafi، \IbnAmir & المدني والشامي\\
\rzlg{سَمَا} & \Nafi، \IbnKathir، \AbuAmr & الحجازيان والبصري\\
\rzlg{حَقّ} & \IbnKathir، \AbuAmr & المكي والبصري\\
\rzlg{نَفَر} & \IbnKathir، \AbuAmr، \IbnAmir & «حق» + ابن عامر\\
\rzlg{حِرْمِيّ} & \Nafi، \IbnKathir & أهل الحرمين\\
\rzlg{حِصْن} & \Nafi{} والكوفيون & «ث» + نافع\\
\bottomrule
\end{tabular}
\end{qaida}

\begin{amthila}[أمثلة لاستعمال الرموز من أبيات الشاطبية الآتية]
\begin{itemize}
\item \rzlg{ث}: «وتظاهرون الظاء خفّف \rzg{ثابتا}» \ila{} الكوفيون: \ayah{تَظَٰهَرُونَ} \src{البقرة: ٨٥} بتخفيف الظاء، والباقون \textbf{تَظَّاهرون} بالتشديد.
\item \rzlg{غ}: «عبادُ برفع الدال في عند \rzg{غلغلا}» \ila{} \AbuAmr{} والكوفيون: \ayah{عِبَٰدُ ٱلرَّحۡمَٰنِ} \src{الزخرف: ١٩}، والباقون \textbf{عِنْدَ الرحمن}.
\item \rzlg{صحبة}: «و\rzg{صحبة} يَصرِفْ فتحُ ضمٍّ» \ila{} \Shuba{} و\Hamza{} و\Kisai: \ayah{مَّن يَصۡرِفۡ عَنۡهُ} \src{الأنعام: ١٦} بفتح الياء وكسر الراء، والباقون \textbf{يُصْرَفْ}.
\item \rzlg{صحاب}: «يُضَلُّ بضم الياء مع فتح ضاده \rzg{صحاب}» \ila{} \Hafs{} و\Hamza{} و\Kisai: \ayah{يُضَلُّ بِهِ ٱلَّذِينَ كَفَرُواْ} \src{التوبة: ٣٧}، والباقون \textbf{يَضِلُّ}.
\item \rzlg{حق}: «و\rzg{حقّ} \rz{Asim}{نصيرٍ} كسرُ واوِ مسوِّمين» \ila{} \IbnKathir{} و\AbuAmr{} و\Asim: \ayah{مُسَوِّمِينَ} \src{آل عمران: ١٢٥} بكسر الواو، والباقون بفتحها.
\item \rzlg{حصن}: «وفي المخلَصين الكل \rzg{حصن}» \ila{} \Nafi{} والكوفيون: \ayah{ٱلۡمُخۡلَصِينَ} بفتح اللام حيث وقعت، والباقون بكسرها.
\end{itemize}
\end{amthila}

\bayt{56}{وَمَهْمَا أَتَتْ مِنْ قَبْلُ أَوْ بَعْدُ كِلْمَةٌ}{فَكُنْ عِنْدَ شَرْطِي وَاقْضِ بِالْوَاوِ فَيْصَلَا}
\maana إذا اجتمع رمز كلمي (من الثمانية) مع رمز حرفي، سواء جاءت الكلمة قبل الحرف أو بعده، فالتزم شرطي: كل رمز يدل على ما وُضع له، ولا يتغير شيء، والواو تفصل بين المسائل كما سبق.

\begin{amthila}
\begin{itemize}
\item «\rzg{وحقّ} \rz{Asim}{نصيرٍ}»: \rzg{حق} = \IbnKathir{} و\AbuAmr، و\rz{Asim}{ن} = \Asim{} \ila{} ثلاثتهم.
\item «\rzg{وعمَّ} \rz{Asim}{علا}»: \rzg{عم} = \Nafi{} و\IbnAmir، و\rz{Asim}{ع} = \Hafs.
\end{itemize}
\end{amthila}

\section{القراءة المسكوت عنها: الأضداد}

\bayt{57}{وَمَا كَانَ ذَا ضِدٍّ فَإِنِّي بِضِدِّهِ}{غَنِيٌّ فَزَاحِمْ بِالذَّكَاءِ لِتَفْضُلَا}
\bayt{58}{كَمَدٍّ وَإِثْبَاتٍ وَفَتْحٍ وَمُدْغَمٍ}{وَهَمْزٍ وَنَقْلٍ وَاخْتِلَاسٍ تَحَصَّلَا}
\bayt{59}{وَجَزْمٍ وَتَذْكِيرٍ وَغَيْبٍ وَخِفَّةٍ}{وَجَمْعٍ وَتَنْوِينٍ وَتَحْرِيكٍ اعْمِلَا}
\begin{mufradat}
\kalima{زاحِم بالذكاء}{اجتهد بفطنتك.}
\kalima{لتفضُلا}{لتكون من الفضلاء.}
\end{mufradat}

\maana[(الأبيات ٥٧–٥٩)] هذه قاعدة عظيمة: الناظم يذكر قراءة من رمز لهم فقط، ويسكت عن الباقين. فإذا كانت القراءة المذكورة لها \textbf{ضد}، فقراءة الباقين تكون بذلك الضد.

\begin{amthila}[مثال]
قال في سورة النساء: «وكوفيهم تساءلون \textbf{مخففًا}». فذكر أن الكوفيين يقرؤون بالتخفيف: \ayah{تَسَآءَلُونَ بِهِۦ} \src{النساء: ١}. وسكت عن الباقين، وضد التخفيف التشديد، فالباقون: \textbf{تَسَّاءَلُونَ}.
\end{amthila}

\begin{qaida}[جدول الأضداد]
\centering
\renewcommand{\arraystretch}{1.4}
\begin{tabular}{c c c}
\toprule
\textbf{إذا ذكر} & \textbf{فضدّه للمسكوت عنهم} & \textbf{هل ينعكس؟}\\
\midrule
المد & القصر & نعم\\
الإثبات & الحذف & نعم\\
الفتح & الإمالة & نعم\\
الإدغام & الإظهار & نعم\\
الهمز & تركه & نعم\\
النقل & تركه (إبقاء الحركة) & نعم\\
الاختلاس & إتمام الحركة & نعم\\
الجزم & الرفع & \textbf{لا} (الرفع ضده النصب)\\
التذكير & التأنيث & نعم\\
الغيب & الخطاب & نعم\\
الخِفّة (التخفيف) & الثِّقَل (التشديد) & نعم\\
الجمع & الإفراد & نعم\\
التنوين & تركه & نعم\\
التحريك & الإسكان & نعم\\
\bottomrule
\end{tabular}

\smallskip\raggedleft
«ينعكس» معناه: إذا ذكر القصر فقراءة الآخرين المد، وإذا ذكر المد فقراءتهم القصر. إلا الجزم: ضده الرفع، أما الرفع فضده النصب لا الجزم.
\end{qaida}

\bayt{60}{وَحَيْثُ جَرَى التَّحْرِيكُ غَيْرَ مُقَيَّدٍ}{هُوَ الْفَتْحُ وَالْإِسْكَانُ آخَاهُ مَنْزِلَا}
\begin{mufradat}
\kalima{غير مقيّد}{لم يذكر نوع الحركة.}
\kalima{آخاه منزلًا}{صار ضده المقابل له.}
\end{mufradat}
\maana إذا قال «حرِّك» ولم يذكر الحركة فالمراد \textbf{الفتح}، وضده الإسكان. أما إذا قيّد الحركة (حرِّك بالضم مثلًا) فالمراد ما قيّده، وضده الإسكان أيضًا.
\begin{amthila}
«معًا \textbf{قدر حرِّك} \rz{IbnAmir}{مِن} \rzg{صِحاب}» \ila{} \IbnDhakwan{} و\Hafs{} و\Hamza{} و\Kisai: \ayah{عَلَى ٱلۡمُوسِعِ قَدَرُهُۥ وَعَلَى ٱلۡمُقۡتِرِ قَدَرُهُۥ} \src{البقرة: ٢٣٦} بفتح الدال، والباقون بإسكانها: \textbf{قَدْرُه}.
\end{amthila}

\bayt{61}{وَآخَيْتُ بَيْنَ النُّونِ وَالْيَا وَفَتْحِهِمْ}{وَكَسْرٍ وَبَيْنَ النَّصْبِ وَالْخَفْضِ مُنْزِلَا}
\maana وجعلت هذه الأزواج أضدادًا أيضًا:
\begin{itemize}
\item \textbf{النون والياء}: «وحيث يشاءُ \textbf{نونُ} \rz{IbnKathir}{دارٍ}» \ila{} \IbnKathir: \ayah{يَتَبَوَّأُ مِنۡهَا حَيۡثُ نَشَآءُ} \src{يوسف: ٥٦} بالنون، والباقون \textbf{يشاء} بالياء.
\item \textbf{الفتح والكسر}: «إنّ الدين \textbf{بالفتح} \rz{Kisai}{رُفِّلا}» \ila{} \Kisai: \ayah{أَنَّ ٱلدِّينَ عِندَ ٱللَّهِ ٱلۡإِسۡلَٰمُ} \src{آل عمران: ١٩} بفتح الهمزة، والباقون \textbf{إنّ} بالكسر.
\item \textbf{النصب والخفض}: «\Hamza{} والأرحامِ \textbf{بالخفض} جمّلا» \ila{} \Hamza: \ayah{وَٱلۡأَرۡحَامِ} \src{النساء: ١} بالجر، والباقون \textbf{والأرحامَ} بالنصب.
\end{itemize}
وفرّق الناظم بين ألقاب البناء (الفتح والكسر) وألقاب الإعراب (النصب والخفض) على اصطلاح البصريين.

\bayt{62}{وَحَيْثُ أَقُولُ الضَّمُّ وَالرَّفْعُ سَاكِتًا}{فَغَيْرُهُمْ بِالْفَتْحِ وَالنَّصْبِ أَقْبَلَا}
\maana إذا قلت «الضم» ولم أذكر ضده، فقراءة الباقين \textbf{بالفتح}. وإذا قلت «الرفع» ولم أذكر ضده، فقراءة الباقين \textbf{بالنصب}.
\begin{amthila}
\begin{itemize}
\item «وفي إذ يرون الياءُ \textbf{بالضم} \rz{IbnAmir}{كُلِّلا}» \ila{} \IbnAmir: \ayah{إِذۡ يُرَوۡنَ ٱلۡعَذَابَ} \src{البقرة: ١٦٥} بضم الياء، والباقون \textbf{يَرَوْن} بالفتح.
\item «وحتى يقولَ \textbf{الرفعُ} في اللام \rz{Nafi}{أوّلا}» \ila{} \Nafi: \ayah{حَتَّىٰ يَقُولُ ٱلرَّسُولُ} \src{البقرة: ٢١٤} بالرفع، والباقون \textbf{يقولَ} بالنصب.
\end{itemize}
أما إذا قيّد الضم بأنه ضد الإسكان أو ضد الكسر، فقراءة الباقين بما قيّده.
\end{amthila}

\bayt{63}{وَفِي الرَّفْعِ وَالتَّذْكِيرِ وَالْغَيْبِ جُمْلَةٌ}{عَلَى لَفْظِهَا أَطْلَقْتُ مَنْ قَيَّدَ الْعُلَا}
\begin{mufradat}
\kalima{أطلقتُ}{ذكرتُ الكلمة دون أن أصفها بالرفع أو التذكير أو الغيب.}
\kalima{من قيّد العُلا}{يا من قيّد المعالي (أي حصّلها) — نداء للطالب.}
\end{mufradat}
\maana قد يذكر الناظم الكلمة بلفظها مطلقة دون قيد، فيُفهم أن المراد: \textbf{الرفع} إن كان الخلاف فيها بين الرفع وغيره، و\textbf{التذكير} إن كان بين التذكير والتأنيث، و\textbf{الغيب} إن كان بين الغيب والخطاب.
\begin{amthila}
\begin{itemize}
\item «وأربعُ أولًا \rzg{صحاب}» \ila{} \Hamza{} و\Kisai{} و\Hafs: \ayah{فَشَهَٰدَةُ أَحَدِهِمۡ أَرۡبَعُ شَهَٰدَٰتِۭ} \src{النور: ٦} بالرفع (ولم يقل «بالرفع» بل أطلقها)، والباقون \textbf{أربعَ} بالنصب. وكلمة «أولًا» هنا معناها: الموضع الأول، وليست رمزًا.
\item «بل يؤثرون \rz{AbuAmr}{حُزْ}» \ila{} \AbuAmr: \ayah{بَلۡ يُؤۡثِرُونَ ٱلۡحَيَوٰةَ ٱلدُّنۡيَا} \src{الأعلى: ١٦} بالغيب (الياء)، والباقون \textbf{تؤثرون} بالخطاب.
\end{itemize}
\end{amthila}

\section{موضع الرموز الكلمية، والتصريح بالأسماء}

\bayt{64}{وَقَبْلَ وَبَعْدَ الْحَرْفِ آتِي بِكُلِّ مَا}{رَمَزْتُ بِهِ فِي الْجَمْعِ إِذْ لَيْسَ مُشْكِلَا}
\begin{mufradat}
\kalima{الحرف}{الكلمة القرآنية.}
\kalima{ما رمزت به في الجمع}{الكلمات الثماني (صحبة، صحاب، عم، سما، حق، نفر، حرمي، حصن).}
\end{mufradat}
\maana الرموز الحرفية (أبج...) \textbf{لا تأتي إلا بعد} الكلمة القرآنية. أما الرموز الكلمية الثمانية فقد آتي بها \textbf{قبل} الكلمة القرآنية أو \textbf{بعدها}؛ لأنها لا تُشكل (فهي كلمات معروفة لا تلتبس).
\begin{amthila}
\begin{itemize}
\item بعد الكلمة: «من يرتدِدْ \rzg{عمّ}» \ila{} \Nafi{} و\IbnAmir: \ayah{مَن يَرۡتَدِدۡ مِنكُمۡ} \src{المائدة: ٥٤} بدالين، والباقون \textbf{يرتدَّ}.
\item قبل الكلمة: «و\rzg{صحبة} يَصرِف» كما سبق.
\end{itemize}
وإذا اجتمع رمز حرفي مع رمز كلمي جاز تقديم الحرف تبعًا للكلمة: «و\rzg{حقّ} \rz{Asim}{نصيرٍ} كسرُ واو مسوّمين».
\end{amthila}

\bayt{65}{وَسَوْفَ أُسَمِّي حَيْثُ يَسْمَحُ نَظْمُهُ}{بِهِ مُوضِحًا جِيدًا مُعَمًّا وَمُخْوَلَا}
\begin{mufradat}
\kalima{أسمّي}{أذكر اسم القارئ صريحًا بدل رمزه.}
\kalima{الجِيد}{العنق.}
\kalima{مُعَمًّا ومُخْوَلًا}{كريم الأعمام والأخوال.}
\end{mufradat}
\maana قد أذكر اسم القارئ صريحًا إذا سمح الوزن، فيكون البيان أوضح، كالعنق الجميل الكريم الأعمام والأخوال (أي واضح الحسن). ويأتي الاسم قبل الكلمة القرآنية أو بعدها.
\begin{amthila}
\begin{itemize}
\item قبل الكلمة: «\Hamza{} والأرحامِ بالخفض جمّلا».
\item في أبوابنا: «عليهم إليهم \Hamza{} ولديهمو» (البيت ١١٠)، و«\Qalun{} بتخييره جلا» (البيت ١١١).
\end{itemize}
\end{amthila}

\bayt{66}{وَمَنْ كَانَ ذَا بَابٍ لَهُ فِيهِ مَذْهَبٌ}{فَلَا بُدَّ أَنْ يُسْمَى فَيُدْرَى وَيُعْقَلَا}
\maana إذا انفرد قارئ أو راوٍ بباب كامل لا يشاركه فيه غيره، فلا بد أن أذكره باسمه الصريح لا برمزه.
\begin{amthila}
\begin{itemize}
\item «ودونك الإدغام الكبير وقطبه \AbuAmr{} البصري» (باب الإدغام الكبير).
\item «و\Hamza{} عند الوقف سهّل همزه» (باب وقف حمزة وهشام على الهمز).
\item «وغلّظ \Warsh{} فتح لام» (باب اللامات).
\end{itemize}
\end{amthila}

\begin{khulasa}[خلاصة الاصطلاحات]
\begin{enumerate}
\item الكلمة القرآنية أولًا، ثم رمز الحرف بعدها (من «أبجد»).
\item الواو تفصل بين المسائل.
\item الرموز الكلمية الثمانية تأتي قبل الكلمة أو بعدها.
\item قد يُصرَّح بالاسم، ولا بد منه إذا انفرد القارئ بباب.
\item يُستغنى باللفظ عن القيد إذا كان اللفظ واضحًا.
\item المسكوت عنهم يُقرأ لهم بضد القراءة المذكورة (جدول الأضداد).
\item التحريك المطلق = الفتح. الضم المطلق ضده الفتح، والرفع المطلق ضده النصب.
\item الكلمة المطلقة بلا قيد = الرفع أو التذكير أو الغيب.
\end{enumerate}
\end{khulasa}
